
\chapter{Überleben II -- Ja, Nein und die Kunst des Feilschens}

\section*{Ja und Nein}
\begin{itemize}[leftmargin=*]
\item \arword{أيوه}{aywa}{Ja (Masri, das gebräuchlichste)}{%D8%A3%D9%8A%D9%88%D9%87}
\item \arword{آه}{aah}{Ja (noch informeller)}{%D8%A2%D9%87}
\item \arword{لأ}{la'}{Nein (Masri, betont)}{%D9%84%D8%A3}
\item \arword{لا، شكراً}{la, shukran}{Nein, danke}{%D9%84%D8%A7%D8%8C%20%D8%B4%D9%83%D8%B1%D8%A7%D9%8B}
\end{itemize}

\begin{lachbox}
Fus'Ha-Nein heißt \textarabic{لا} (\textit{la}) -- kurz, unschuldig,
harmlos. Masri-Nein heißt \textarabic{لأ} (\textit{la'}) -- mit einem
kleinen Kehlkopf-Stopp am Ende, der Entschlossenheit signalisiert. Julia,
lern dieses Nein. Du wirst es auf dem Basar brauchen.
\end{lachbox}

\section*{Feilschen -- eine Nationalsportart}
\begin{itemize}[leftmargin=*]
\item \arword{بكام ده؟}{bikaam da?}{Wie viel kostet das?}{%D8%A8%D9%83%D8%A7%D9%85%20%D8%AF%D9%87%D8%9F}
\item \arword{غالي أوي!}{gháali awi!}{Viel zu teuer!}{%D8%BA%D8%A7%D9%84%D9%8A%20%D8%A3%D9%88%D9%8A%21}
\item \arword{كتير أوي}{kitír awi}{sehr/zu viel}{%D9%83%D8%AA%D9%8A%D8%B1%20%D8%A3%D9%88%D9%8A}
\item \arword{ممكن تنزل السعر؟}{mumkin tinzil is-si3r?}{Kannst du mit dem Preis runtergehen?}{%D9%85%D9%85%D9%83%D9%86%20%D8%AA%D9%86%D8%B2%D9%84%20%D8%A7%D9%84%D8%B3%D8%B9%D8%B1%D8%9F}
\item \arword{ماشي}{máashi}{Okay, abgemacht}{%D9%85%D8%A7%D8%B4%D9%8A}
\item \arword{والله؟}{walláhi?}{Wirklich?! (bei Gott, Ausdruck des Erstaunens)}{%D9%88%D8%A7%D9%84%D9%84%D9%87%D8%9F}
\end{itemize}

\begin{ammahmoudbox}
,,Der erste Preis, den ein Verkäufer nennt, ist wie das erste Angebot bei
einem Flohmarkt in Bayern: eine Einladung, nicht eine Tatsache. Wenn du
nicht handelst, ist das für uns fast unhöflich -- es nimmt uns den Spaß.``
\end{ammahmoudbox}

\begin{tippbox}
Faustregel für Touristen-Basare: biete zuerst \emph{die Hälfte} des
genannten Preises, lächle dabei, und lass dich langsam Richtung Mitte
verhandeln. Bei \textarabic{غالي أوي!} (,,viel zu teuer!``) darfst du
theatralisch die Augenbrauen hochziehen -- das gehört dazu.
\end{tippbox}

\begin{checkpointbox}
Rollenspiel-Vorbereitung: du wirst \textarabic{بكام ده؟},
\textarabic{غالي أوي!} und \textarabic{ماشي} in Teil 8 in einem echten
Dialog mit Am Mahmoud brauchen. Sprich sie schon jetzt laut.
\end{checkpointbox}
